\problem{P21}{两种整数步长的更新过程}{递推与越界量}
独立增量 $X_i$ 取值 $1,2$，概率分别为 $1-q,q$，其中 $0\le q\le1$。对整数 $m\ge1$，令 $T_m=\min\{n:\sum_{i=1}^nX_i\ge m\}$，越界量为 $O_m=\sum_{i=1}^{T_m}X_i-m$。求 $\E T_m$ 和 $O_m$ 的分布。
\solution 记 $e_m=\E T_m$，并置 $e_0=e_{-1}=0$。按第一步条件化，
\[
e_m=1+(1-q)e_{m-1}+qe_{m-2},\quad m\ge1.
\]
令 $d_m=e_m-e_{m-1}$，则 $d_1=1$，且 $d_m=1-qd_{m-1}$。由等比数列求和，
\[
d_m=\frac{1-(-q)^m}{1+q},\qquad
\boxed{e_m=\frac m{1+q}+\frac{q[1-(-q)^m]}{(1+q)^2}.}
\]
每步至少前进 $1$，故 $T_m\le m$。写成有限指示变量和，利用 $\{T_m\ge i\}$ 只依赖前 $i-1$ 个增量，得到
\[
\E\sum_{i=1}^{T_m}X_i
=\sum_{i=1}^m\E[X_i\ind\{T_m\ge i\}]
=(1+q)e_m.
\]
越界量仅能取 $0,1$，所以
\[
\boxed{\Prob(O_m=1)=\E O_m
=\frac{q[1-(-q)^m]}{1+q},\quad
\Prob(O_m=0)=1-\Prob(O_m=1).}
\]
$q=0$ 时恰需 $m$ 步；$q=1$ 时恰需 $\lceil m/2\rceil$ 步，越界量由 $m$ 的奇偶性决定，均与公式相符。
\insight 整数步长留下了交替项 $(-q)^m$。这项反映格点结构，不能用仅含平均步长的近似替代。
